% !TEX root = ../main.tex
% =====================================================================
% 8. FINDING 4 -- capacity-safe is not optimisation-safe. Budget 0.75.
%
% A negative result. It is only publishable because the mechanism is
% identified, so the mechanism paragraph is the important one, not the
% table.
% =====================================================================

\section{Capacity-Safety Is Not Optimisation-Safety}
\label{sec:safety}

\begin{finding}
\label{find:safety}
Five of our seven heads provably reduce to the conjunctive baseline for one setting of their
parameters. That property is routinely used to argue a design ``cannot hurt''. One such head is
nevertheless the worst we tested, by a wide margin, and the reason is that no term in the objective
drives the model back to the safe branch.
\end{finding}

The fourth column of \Cref{tab:heads} is an argument that appears often in the attention and MIL
literature: if a new head contains the established one as a special case, then in the worst case the
optimiser recovers the established one, so adding the head is free. \Cref{tab:pooling} tests that
argument on eight heads averaged over every encoder they were paired with.

\begin{table}[t]
\caption{\webtraffic{} accuracy, AOPCR and \ndcg{} for each pooling head, averaged over the $n$
encoders it was paired with. ``Reduces'' repeats the fourth column of \Cref{tab:heads}. The
$\kappa$-ablation configurations and the alternative-recipe baseline are excluded, so that no
pairing is counted twice. Groups with small $n$ are not comparable with groups averaged over ten or
more pairings; the full grid is in \Cref{app:ablations}.}
\label{tab:pooling}
\begin{center}
\footnotesize
\begin{tabular}{lccccc}
\toprule
\textbf{Pooling head} & \textbf{Reduces} & \textbf{Mean acc.} & \textbf{Mean AOPCR} &
\textbf{Mean \ndcg{}} & \textbf{$n$} \\
\midrule
Class-wise conjunctive      & yes & \textbf{0.921} & 2.322 & 0.692 & 10 \\
Per-class gated attention   & \textbf{no}  & 0.918 & 1.869 & \textbf{0.733} & \phantom{0}4 \\
Adaptive class-wise         & yes & 0.912 & 2.210 & 0.732 & \phantom{0}9 \\
Softmax conjunctive         & yes & 0.907 & 1.690 & 0.712 & \phantom{0}9 \\
Top-$k$ conjunctive         & yes & 0.904 & 2.267 & 0.712 & 16 \\
Attention-max               & yes & 0.838 & 1.758 & 0.645 & \phantom{0}8 \\
\textbf{Dual-stream conjunctive} & \textbf{yes} & \textbf{0.744} & 2.007 & 0.636 & \phantom{0}4 \\
\bottomrule
\end{tabular}
\end{center}
\end{table}

\paragraph{The counterexample.}
Dual-stream conjunctive pooling blends the class-wise conjunctive mean with a critical-instance
stream through a single learnable scalar $\lambda = \sig(\theta)$, initialised at $0.05$. At
$\lambda \to 0$ it is the baseline exactly. It is nonetheless the weakest head in the study: mean
accuracy $0.744$ against $0.90$ and above for the three strongest heads, with its four pairings at
$0.860$, $0.824$, $0.736$ and $0.556$. A head that provably cannot be worse than $0.921$ scored
$0.556$ on one pairing.

\paragraph{The mechanism.}
The property is about \emph{representational capacity}; the failure is one of \emph{optimisation}.
Nothing in \Cref{eq:loss} rewards small $\lambda$. If the critical-instance query learns a poor
embedding early in training, the gradient through the second stream is not informative, but there is
also no pressure returning $\lambda$ towards zero -- the model sits in a region where the blend is
harmful and the safe branch, though reachable in principle, is never searched for. The model
\emph{can} represent the baseline; it never finds it. This is not specific to the dual stream: any
``contains the baseline as a special case'' argument has the same gap between what the parameter
space allows and what the optimiser visits.

\paragraph{The contrast that supports the reading.}
Two comparisons inside \Cref{tab:pooling} sharpen it. Top-$k$ pooling has the same reduction
property but implements it through a \emph{hyperparameter} rather than a learned weight: at
$\kappa = 1$ the baseline is recovered by construction, not by optimisation, and Top-$k$ never
collapses. Conversely, per-class gated attention is the only head with \emph{no} reduction at all,
and its behaviour is the mirror image -- a respectable mean of $0.918$ and the best mean \ndcg{}, but
never the best partner for any encoder. It has no known-good configuration to fall back to, so it
neither wins nor fails outright.

\paragraph{What a designer should take from this.}
A reduction property is worth having, but it should be built where the optimiser cannot lose it:
prefer a fixed hyperparameter over a learned blend, initialise a learned blend at the safe value so
that leaving it is a deliberate move, or add an explicit penalty holding the blend near the safe
branch early in training. We initialised $\lambda$ at $0.05$ rather than $0$, which we now regard as
the error; we did not have the compute to re-run the alternative, so this is a diagnosis with a
proposed remedy rather than a demonstrated fix.
